\subsection*{S0. Unit of input}

JSON reaches the tool in more than one shape, and the shapes disagree about
what counts as one document. The distinction matters because inference is over
the set of values at a path: a root field is observed once per document, so
how many documents there are determines how much can be inferred.

Two groups of input form must be separated. The first group changes what
counts as a document and therefore needs a decision:

\begin{itemize}
  \item \textbf{Single document} --- one JSON value.
  \item \textbf{Top-level array} --- \texttt{[\{...\},\{...\}]}. Ambiguous:
        either a corpus of several documents, or one document whose root
        happens to be an array.
  \item \textbf{JSON Lines / NDJSON} --- one JSON value per line.
        Unambiguously a corpus.
  \item \textbf{Concatenated JSON} --- \texttt{\{...\}\{...\}} with no
        delimiter. A corpus, but requires a streaming parser.
  \item \textbf{Wrapped corpus} --- \texttt{\{"meta":\{...\},"records":[...]\}},
        where the documents live at a nested path rather than at the root.
\end{itemize}

The second group is plumbing: a directory of files, several files on the
command line, standard input, compression, character encoding. These change
how bytes reach the parser and nothing about inference, so they carry no
design decision and can be added at any point.

\options
\begin{enumerate}[label=\textbf{(\alph*)}]
  \item Single document only. Inference from one JSON value.
  \item Corpus only. The input is always many documents.
  \item Both, with a top-level array always read as a corpus. \selected
  \item Both, disambiguated by an explicit \texttt{-{}-corpus} /
        \texttt{-{}-single} flag.
\end{enumerate}

\decision{Accept a single document and a top-level array. A top-level array is
always read as a \emph{corpus of its elements}, never as one document whose
root is an array. JSON Lines, concatenated JSON and wrapped corpora are
deferred; they extend the input reader without disturbing any other decision
in this document.}

Given a single document, the corpus is that one document:

\begin{lstlisting}[style=json]
{ "a": 1, "b": { "c": 2 } }

  -> corpus of 1 document
\end{lstlisting}

Given a top-level array, the corpus is its elements:

\begin{lstlisting}[style=json]
[ { "a": 1 }, { "a": 2 } ]

  -> corpus of 2 documents
  -> NOT one document whose root is an array
\end{lstlisting}

\rationale{Reading a top-level array as a corpus removes the ambiguity without
a flag, and matches how JSON corpora are actually distributed. It also implies
that the root of a document is always an object, which is why S12 no longer
needs to exist as a separate scenario.}

\limitation{A document whose root is genuinely an array cannot be expressed;
it must be wrapped, as \texttt{\{"items": [...]\}}. Separately, a single
document cannot support optionality inference: every field present in it is
total and no field is ever partial, so single-document input yields a strictly
weaker schema than corpus input. Arrays \emph{inside} a single document still
supply many observations, so widening and optionality behave normally within
them.}
